% Paper 8 Figure 6: Conditional analytical illustration of the single-edge residual.

\begin{figure}[htbp]
\centering
\begin{tikzpicture}
\begin{axis}[
  harbor curve,
  width=0.7\linewidth,
  height=6cm,
  title style={font=\sffamily\bfseries\small, yshift=4pt},
  title={Analytical Illustration: Assume $R^K_{\mathrm{eff}}(e)=1/3$},
  xlabel style={font=\sffamily\small},
  xlabel={Single-Edge Lie Magnitude $|s|$},
  ylabel style={font=\sffamily\small},
  ylabel={Completion Residual $r$},
  xmin=0, xmax=5.2,
  ymin=0, ymax=4.6,
  grid=both,
  grid style={line width=0.35pt, draw=ptGrid},
  major grid style={line width=0.55pt, draw=ptGrid!75!black},
  legend pos=north west,
  legend cell align={left},
  legend style={font=\sffamily\scriptsize, draw=ptDark, line width=0.8pt, fill=white},
  tick label style={font=\sffamily\scriptsize},
  axis line style={draw=ptDark, line width=0.9pt},
]
\addplot[
  domain=0:5,
  samples=2,
  color=ptBlue,
  line width=1.4pt
] {x * sqrt(2/3)};
\addlegendentry{$r=|s|\sqrt{2/3}$}
\end{axis}
\end{tikzpicture}
\caption{Conditional analytical illustration of the single-edge formula $r=|s|\sqrt{1-R^K_{\mathrm{eff}}(e)}$. Assuming effective resistance $R^K_{\mathrm{eff}}(e)=1/3$ in the unit-weight visible graph gives $r=|s|\sqrt{2/3}$. Every point on this curve is calculated from that assumed value and the formula.}
\label{fig:matrix}
\end{figure}
